\section{Conditional reconstruction from finitely many noisy windows}
\label{sec:stable}
Exact analytic continuation and a finite-dimensional coordinate inverse
are different statements. We now give a convergent regularization for
the whole table. Its hypotheses include quantitative analytic and
asymmetry bounds; its observations include noisy onset localization.
No boundary image, contact registration, or exact calibration limit is
an input. The rate below is a conditional upper bound, not an assertion
of minimax optimality or uniform conditioning of the jet recursion.

\subsection{A priori class and comparison metric}
Fix the labelled full-rank network and a rooted spanning tree. A change
of representative copies makes its tree labels zero. Choose two
independent fundamental cycle labels and put
$\mathsf K=[k'_{e_1}\ k'_{e_2}]$. The following bounds define the
admissible class $\mathcal A$.
The number of obstacles, their diameters, and the lattice matrix and its
inverse are bounded. The gap, curvature, free area, global obstacle separation, channel
isolation and clearance have fixed positive lower bounds and fixed upper bounds
where appropriate. Each onset belongs to a known bounded bracket; the
experiment can query a fixed collar beyond the upper bracket endpoint.
The brackets concern the selected local event, not an unknown competing
itinerary. There is a common bound $S$ on the absolute endpoint marks.

In any centered orthonormal frame each obstacle support function extends
holomorphically to $|\operatorname{Im}\theta|<\rho$ and is bounded there
by $M$. Local contact graphs extend to $|y|<R$ with a common bound.
We work with a closed class satisfying these bounds and the preceding
strict geometric margins. Equivalently one may take a closed subclass
of a sufficiently small uniformly analytic neighborhood with those
margins. All constants below are allowed to depend on this class. The
holomorphic graph bound can also be obtained, with smaller $R$, from
uniform support bounds and positive curvature by the analytic
inverse-function theorem.

Write $p_C$ for the support function after subtracting the Steiner point.
For some fixed $\eta>0$ require
\begin{align}
 \|p_C(\cdot+\phi)-p_C\|_{L^2(0,2\pi)}
   &\ge\eta\,\operatorname{dist}(\phi,2\pi\mathbb Z),
                                      \label{eq:rotationmargin}\\
 \|p_C(-\cdot+\phi)-p_C\|_{L^2(0,2\pi)}
   &\ge\eta,\qquad \phi\in\R.          \label{eq:reflectionmargin}
\end{align}
Every asymmetric analytic body satisfies these inequalities with some
positive $\eta$. For rotations close to zero this follows from
$\|p_C'\|_2>0$; away from zero and for reflections it follows from
compactness and absence of symmetry. On a sufficiently small compact
neighborhood of such a body the same smaller margin works. Thus these
priors leave infinite-dimensional physical classes, including small
uniformly analytic perturbations of \eqref{eq:asymmetricsupport}.
They do not specify a finite-dimensional model of the boundary.

For two tables in the tree representative gauge, let
$\operatorname{dist}_{\rm tab}$ be the infimum, over one common Euclidean
isometry, of the maximum of the $C^2$ distances between their labelled
support functions and the norm of the difference of their marked
lattice matrices. Translations of representatives are included in the
support functions. The orthogonal part of the isometry acts on the
lattice matrix. A bounded root placement may be fixed when taking
compactness. Montel compactness on a narrower strip and the geometric
margins make $\mathcal A$ compact in this metric, after reducing the
strip once if necessary.

\begin{theorem}[Finite-data regularization of the entire table]
\label{thm:stable}
For the class $\mathcal A$ there are constants $C,B>0$, $0<q<1$,
$0<\vartheta<1$, and, for every $K\ge3$, finite constants
$A_K>0$ and $\epsilon_K>0$ with the following property. A finite
adaptive experiment uses only masses and the two endpoint moments
$M_b,D_b$ at positive absolute times. If all its scalar means have
error at most $\epsilon\le\epsilon_K$, its admissible reconstruction
$\widehat{\mathcal T}$ satisfies
\begin{equation}\label{eq:tablestability}
 \operatorname{dist}_{\rm tab}(\widehat{\mathcal T},\mathcal T)
 \le C\left(A_K\epsilon^{1/(n_K+1)}+Bq^K\right)^{\vartheta},
 \qquad n_K=\lceil K/2\rceil+1.
\end{equation}
The experiment observes each channel with its own reversal orbit. It
supplies neither relative edge signs nor contact registrations, and it
estimates the onset, gap, separate curvatures and area rather than
receiving them exactly. It uses
$O(|E|(\epsilon^{-1/2}+n_K+1))$ scalar mean channels.
There is a choice $K=K(\epsilon)\to\infty$ for which the right side
of \eqref{eq:tablestability} tends to zero. No uniform growth bound on
$A_K$ is required or asserted.

For independent fresh phase preparations and error probability
$0<\delta<1/4$, the required simultaneous mean accuracy is achieved
with at most
\begin{equation}\label{eq:fullsamplebudget}
 C|E|(\epsilon^{-1/2}+n_K+1)\epsilon^{-2}
 \log\!\left(\frac{C|E|(\epsilon^{-1/2}+n_K+1)}{\delta}\right)
\end{equation}
preparations, including failures. The guarantee then has probability
at least $1-\delta$. This is an upper bound for the specified marked
experiment, not a lower bound for a count-only or full-trajectory sensor.
\end{theorem}

\subsection{Locating the onset and estimating finite jets}
The common analytic bounds, the positive Hessian, and the stationary
and Morse constructions in Lemma~\ref{lem:flux} give a common right
analytic collar for $P_b,M_b,D_b$. Their derivatives of every fixed
order are uniformly bounded on a smaller collar. To see the uniformity,
apply the analytic implicit-function theorem on the bounded complex
graph neighborhoods, retaining a smaller common neighborhood where the
stationary derivative stays away from zero. The positive Hessians
admit parameter-uniform Morse charts there. Integration on the fixed
Morse disk then preserves the uniform analytic bounds. In particular,
for fixed positive $c,C,d_*$,
\begin{equation}\label{eq:onsetbounds}
 P_b(T_0+d)=\alpha d^2+O(d^3),\quad
 c d^2\le P_b(T_0+d)\le C d^2\quad(0<d<d_*).
\end{equation}
Before onset this probability is zero. The right analytic extension is
used for Taylor estimates only; it is not claimed to equal the physical
probability at negative offsets.

On each known bracket query $P_0$ on a grid of spacing
$\sqrt\epsilon$, adding a fixed collar beyond the bracket endpoint.
Let $\tau$ be the first grid point whose reported mean exceeds
$8\epsilon$. If every error is at most $\epsilon$, no point at or
below $T_0$ can trigger. A point within $C_0\sqrt\epsilon$ after
$T_0$ has probability greater than $9\epsilon$ by
\eqref{eq:onsetbounds}. Thus, for small enough $\epsilon$,
\begin{equation}\label{eq:onsetestimate}
             0<\tau-T_0\le C_1\sqrt\epsilon.
\end{equation}
This argument uses no global monotonicity of the three-impact event.
Its first guaranteed crossing lies inside the uniform local collar.
The grid contains $O(\epsilon^{-1/2})$ points.

Put $n=n_K$ and $h=\epsilon^{1/(n+1)}$. At times
$\tau+h,\ldots,\tau+(n+1)h$ estimate each of the six means
$P_b,M_b,D_b$, $b=0,1$, and interpolate a polynomial of degree $n$
in $x=T-\tau$. For every fixed $n$, Taylor expansion at $\tau$ and
the inverse of the fixed nodal matrix $(i^j)_{1\le i\le n+1,0\le j\le n}$
give, for each estimated coefficient of order $0\le j\le n$,
\begin{equation}\label{eq:coefficientnoise}
 |\widehat a_j-a_j(T_0)|
 \le C_n\{\sqrt\epsilon+h+\epsilon h^{-n}\}
 \le C_n'\epsilon^{1/(n+1)}.
\end{equation}
Indeed the Taylor remainder at node $ih$ is $O_n(h^{n+1})$,
so its coefficient error is $O_n(h^{n+1-j})$; the observation error
is $O_n(\epsilon h^{-j})$. Translation from $\tau$ to $T_0$
changes every fixed coefficient by $O_n(|\tau-T_0|)$, since the
right analytic derivatives are uniformly bounded. This also explains
why an unknown onset is not silently inserted into the interpolation.
The threshold $\epsilon_K$ ensures that all these windows lie in the
common collar.

The coefficient of $d^2$ in $P_b$ is $\alpha$, and that of $d^3$
in $D_b$ is $\alpha\beta_b$. The mass coefficients through
$d^{\lfloor K/2\rfloor+1}$ and the first-moment coefficients through
$d^{\lfloor(K-1)/2\rfloor+2}$ suffice for the $K$-jet inverse.
Their maximum degree, together with three for calibration, is precisely
$n_K$. Positive leading margins, \eqref{eq:calibrationmain}, and the
finite triangular recursion give
\begin{equation}\label{eq:jetnoise}
 |\widehat g-g|+|\widehat A-A|
 +\sum_{b=0}^1\left(|\widehat\kappa_b-\kappa_b|
          +\sum_{r=3}^K|\widehat q_{b,r}-q_{b,r}|\right)
 \le B_K\epsilon^{1/(n_K+1)}.
\end{equation}
The constants can be obtained recursively by bounding the derivatives
of the explicit finite block inverse and of its finite-jet remainders
on the prior box. This requires no inverse estimate at infinite order.

\subsection{Propagation from contact arcs to closed boundaries}
We include the continuation estimate to distinguish its loss from the
finite-jet conditioning. It is the elementary two-constants mechanism
behind conditional analytic continuation; see also \cite{Trefethen}
for the dependence of such losses on the continuation geometry.

\begin{lemma}[Conditional continuation on the periodic strip]
\label{lem:continuationbound}
Let $p,p'$ be $2\pi$-periodic holomorphic functions bounded by $M$
on $|\operatorname{Im}\theta|<\rho$. If their difference is bounded
by $e\le2M$ on a real interval of fixed positive length, then
\begin{equation}\label{eq:stripstability}
 \|p-p'\|_{C^2(\R/2\pi\mathbb Z)}\le C e^{\vartheta}
\end{equation}
for some $C>0$ and $0<\vartheta<1$ depending only on $M,\rho$
and that length.
\end{lemma}
\begin{proof}
Choose a closed shorter interval $I$ within the observed interval and
a narrower strip of width $\rho_0<\rho$. Regard this strip modulo
$2\pi$ as a finite cylinder, slit along $I$. Let $\omega$ be the
harmonic function with boundary values one on both sides of the slit
and zero on the two outer circles. These piecewise smooth boundaries
are regular for the Dirichlet problem; the endpoints of the slit have
barriers. The maximum principle gives $0<\omega<1$ in the connected
slit cylinder. On a still narrower closed cylinder, extending by one
on $I$, its minimum is a positive number $\vartheta$.

Apply the subharmonic maximum principle to $\log|p-p'|$, using the
bound $\log e$ on the slit and $\log(2M)$ elsewhere. Zeros cause no
problem by regularization and a decreasing limit. This gives
$|p-p'|\le(2M)^{1-\omega}e^\omega$ and hence a uniform
$C_0e^{\vartheta}$ bound on the narrower cylinder. Cauchy's estimates
on fixed disks in that cylinder give the first and second derivative
bounds on the real circle. The case $e=0$ follows by the identity
theorem or by taking a limit. The constants are uniform under translating
$I$ around the cylinder.
\end{proof}

For two admissible tables compatible with the same finite observations,
\eqref{eq:jetnoise} bounds the differences of their contact jets, after
choosing the compared representative separately on each edge. Uniform
holomorphic graph bounds and Taylor's formula on nested disks give
\begin{equation}\label{eq:localanalyticerror}
 \|\psi_b-\psi'_b\|_{C^2([-r,r])}
       \le A_K\epsilon^{1/(n_K+1)}+Bq^K=:e_K,
                    \qquad 0<r<R,
\end{equation}
with $B,q$ independent of $K$. One may first bound the tail on an
intermediate complex disk and use Cauchy's estimates on a smaller disk;
this absorbs the derivative factors without changing the geometric
tail form. The finite polynomial part is absorbed into $A_K$.

Positive curvature and the uniform $C^3$ bounds identify a common
normal-angle interval inside each of these two graph arcs. The map
from the graph to its support function there is uniformly locally
Lipschitz in $C^2$: the normal angle is obtained by inverting
$\psi_b'$, whose derivative is bounded below, and the point is
$(-\psi_b(y),y)$ or its translated opposite-frame version.
The support value is its scalar product with the corresponding normal.
The inverse-function derivative formulas, with the stated bounds,
therefore give an $O(e_K)$ support error on a smaller fixed interval.
Lemma~\ref{lem:continuationbound} now gives $C e_K^{\vartheta}$
error for the entire recovered boundary in each local contact frame.
The gap error in \eqref{eq:jetnoise} is absorbed in the same bound.
Changing from Steiner-centered frames to these contact frames enlarges
the complex support bound only by a known multiple of the diameter and
$\cosh\rho$. Thus each whole pair realization is close, not merely
its two germs.

\subsection{Stable synchronization and the lattice}
Put $\Delta=C e_K^{\vartheta}$ and suppose it is smaller than a
fixed fraction of $\eta$. The Steiner point is a bounded linear
functional of the support function, so centering changes an error of
size $\Delta$ by at most a fixed multiple. Suppose two ways of
matching a source body to an already placed image both fit to this
accuracy. Their relative orthogonal transformation moves the centered
support function by $O(\Delta)$ in $L^2$. Equation
\eqref{eq:reflectionmargin} excludes the wrong reflection component;
\eqref{eq:rotationmargin} bounds the difference of their rotation
angles by $C\Delta/\eta$. The bounded $C^3$ support norms then
bound their action on the complete pair in $C^2$ by $C\Delta/\eta$.
Their translation error is controlled by the Steiner points.

More explicitly, when comparing two physical candidate tables, each
edge pair is already known to be $O(\Delta)$-close in its own frame.
Its matching to the placed source in either table therefore gives two
approximate matches of the same source image. The preceding margin
argument aligns them, with no independent choice of a sign. Induction
along the finite spanning tree bounds the differences of all placed
representative bodies by $C_{G,\eta}\Delta$. On a non-tree edge the
same argument controls its target image. Taking Steiner-point
differences in \eqref{eq:unregisteredcycle} gives the same bound for
$d_e$. Consequently
\begin{equation}\label{eq:latticestability}
 \|L-L'\|\le C_{G,\eta}\|\mathsf K^{-1}\|\Delta.
\end{equation}
All remaining cycle equations are consistency tests. This argument
compares admissible physical tables; it does not declare an arbitrary
noisy curve to be an exactly congruent obstacle. For larger $e_K$ the
bounded diameter of the prior class supplies the estimate after
increasing $C$. Equations \eqref{eq:localanalyticerror}--\eqref{eq:latticestability}
prove \eqref{eq:tablestability} for any compatible reconstruction.

\subsection{Regularization and the statistical experiment}
For clarity a compatible reconstruction can be defined without assuming
an efficient nonlinear solver. On the compact admissible class minimize
the maximum discrepancy from the finitely many reported means, allowing
one sign choice per edge, common to that edge's windows. Choose an
approximate minimizer from a fixed countable dense family with discrepancy
within $\epsilon$ of the infimum. Such a choice may be made Borel by
an ordered first-eligible-point rule. The true table has discrepancy at
most $\epsilon$ on the mean-accuracy event. The selected table and
the truth therefore have discrepancies bounded by a fixed multiple of
$\epsilon$, and the preceding proof applies with rescaled constants.
The first-stage threshold $8\epsilon$ exceeds every compatible
candidate's allowed discrepancy. The onset argument therefore applies
to each candidate as well as to the truth, with adjusted constants:
its onset is below the first crossing and within $C\sqrt\epsilon$
of that crossing. For the two-stage design the first-stage
data are retained in this minimization, not discarded. The mean functions are continuous on this class. Indeed the uniform
obstacle separation bounds the number of collisions in every queried
bounded time interval. Outside the zero-volume sets of grazing rays,
patch-boundary impacts and terminal-time impacts, the finite collision
record and its arclength marks vary continuously with the table.
A fixed cell-coordinate parametrization and the positive area bound
then allow dominated convergence, including the failure indicator.
This continuity and compactness justify the dense-family construction.
The estimator is an existence and consistency result, not a polynomial-time
algorithm for searching the admissible class.

Choose a decreasing sequence $t_K\downarrow0$ such that
$t_K\le\epsilon_K$ and $A_K t_K^{1/(n_K+1)}\le q^K$.
For $t_{K+1}<\epsilon\le t_K$ use order $K$ (and take a fixed
initial order for larger errors). Then $K\to\infty$ and the right
side of \eqref{eq:tablestability} tends to zero. This gives a
convergent, explicitly regularized finite-data inverse even if the
finite-jet constants deteriorate rapidly. It does not impose an
unproved convergence estimate on a formal count-only inverse.

Finally each queried random variable is bounded in absolute value by
$\max(1,2S,4S^2)$, with zero mark on failure. Independent sample means
therefore have tails $2\exp(-cN\epsilon^2)$. The first-stage grid is
deterministic. Conditional on its record, the second-stage windows are
fixed and fresh independent preparations have the same bound. A union
bound over both stages and all edges gives simultaneous accuracy with
$C\epsilon^{-2}\log(CJ/\delta)$ preparations per channel, where
$J\le C|E|(\epsilon^{-1/2}+n_K+1)$. This proves
\eqref{eq:fullsamplebudget} and completes Theorem~\ref{thm:stable}, hence
Theorem~\ref{thm:intrinsic}(v).

The analytic strip and symmetry margins are substantive conditions of
this result. Without a uniform analytic prior the identity theorem alone
supplies no noisy continuation estimate; near symmetric obstacles
unregistered synchronization can be poorly conditioned. We assert
neither an optimal exponent for the full inverse nor a solution of the
exact analytic count-only quotient problem. The fully registered
symmetric-table theorem and the separate fixed-dimensional experiments
remain valid in their stated information categories.
